\exercisesection{5}

\begin{enumerate}
\def\labelenumi{\arabic{enumi}.}
\item
  给出一个自同态 \(u\) ，它定义在 \(\mathbf{Z}\) -模 \(\mathbf{M} = \mathbf{Z}^{(\mathbf{N})}\) 上，使得：当 \(p = (0)\) 是 \(\mathbf{Z}\) 的素理想时，\(u_{p}\) 是 \(M_p\) 的自同构，但不存在 \(f \neq 0\) （\(\mathbf{Z}\) 中的）使 \(u_f\) 是 \(M_f\) 的满射自同态。
\item
  设 K 是一个域，B 是以无限未定元系 \(X_{i}\) （ \(i \in N\) ）为未定元、系数在 K 中的多项式组成的环；设 b 是 B 的由诸乘积 \(\mathbf{X}_i\mathbf{X}_j\) （ \(i \neq j\) ）生成的理想，A 是商环 B/b，\(\xi_i\) 是 \(\mathbf{X}_i\) 模 b 所在的类（ \(i \in \mathbb{N}\) ），p 是 A 的由诸 \(\xi_i\) （其指标 \(i \geqslant 1\) ）生成的理想（它是素的），u 是典范映射 A → A/p。证明 \(u_p: \mathbf{A}_p \to (\mathbf{A}/\mathfrak{p})_p\) 是双射，但不存在 \(f \in \mathbf{A} - \mathfrak{p}\) 使得 \(u_f\) 是双射。
\item
  设 A 是一个环，P 是一个有限生成投射 A-模。证明 A 同构于环的一个有限乘积 \(\prod_{i\in I}A_{i}\) ，且 P 同构于一个乘积 \(\prod_{i\in I}P_{i}\) ，其中 \(P_{i}\) 是 \(A_{i}\) 上的秩为 \(n_{i}\) 的投射模，而诸 \(n_{i}\) 互不相同（利用定理 1 及 \(\operatorname{Spec}(A)\) 是拟紧的这一事实）。
\item
  设 A 是一个（交换）环，B 是一个包含 A 的（不必交换的）环。假设左 A-模 B 是投射的且有限生成的。证明 A 是 B 的直因子。（归结到 A 是局部环的情形，并利用 § 3 第 2 节命题 5 的推论 1。）
\item
  设 A 是一个约化环，M 是一个有限生成 A-模。假设存在一个整数 n \textgreater{} 0，使得对每个同态 \(\phi\) （由 A 到一个域 \(K_{\phi}\) ），都有 \([M \otimes_{A} K_{\phi}: K_{\phi}] = n\) 。证明 M 是秩为 n 的投射模。（归结到 A 是局部环的情形，并利用 §2 第 2 节命题 7。）
\end{enumerate}

¶ 6. 设 A 是一个整环，K 是它的分式域，M 是一个 A-模。证明下列性质等价：

\((\alpha)\) M 是一个投射 A-模，且 \([\mathbf{M} \otimes_{\mathbf{A}} \mathbf{K} : \mathbf{K}]\) 是有限的。

(β) M 是有限生成投射 A-模。

\((\gamma)\) M 是有限生成的，并且对 A 的每个极大理想 m，\(\mathbf{A}_{m}\) -模 \(\mathbf{M}_{m}\) 是自由的。

（为证明 \((\alpha)\) 蕴含 \((\beta)\) ，利用《代数》第 II 章 § 5 第 5 节命题 9。为证明 \((\gamma)\) 蕴含 \((\alpha)\) ，首先注意 M 是无扭的，然后断言：对 A 的每个素理想 \(p\) ，\(M_{p}\) 是有常数秩的自由 \(A_{p}\) -模。）

\begin{enumerate}
\def\labelenumi{\arabic{enumi}.}
\setcounter{enumi}{6}
\tightlist
\item
  设 A 是一个绝对平坦的（交换）环（参看 § 4 习题 16），X 是它的谱。
\end{enumerate}

\begin{enumerate}
\def\labelenumi{(\alph{enumi})}
\item
  设 F 是 X 的一个闭子集，a 是 A 的相应理想（出处同上）。证明 M = A/a 是一个单生成的平坦 A-模，使得 \(M_{m}\) 是自由 \(A_{m}\) -模（对所有 \(m \in X\) ），但仅当 F 是开的时 M 才是投射的。
\item
  设 E 是诸 A/m（ \(m \in X\) ）的直和组成的模。证明 E 是一个平坦 A-模，使得 \(E_{m}\) 是秩为 1 的自由 \(A_{m}\) -模（对所有 \(m \in X\) ）；为使 E 是一个投射 A-模，必须且只需 X 是有限的（注意：若 A/m 是一个投射 A-模，则 m 是有限生成的）。
\end{enumerate}

\begin{enumerate}
\def\labelenumi{\arabic{enumi}.}
\setcounter{enumi}{7}
\tightlist
\item
  设 A 是一个环，M 是一个秩为 n 的投射 A-模。证明存在一个有限生成的（交换）A-代数 B，使得 A-模
\end{enumerate}

B 是忠实平坦的，且 \(\mathbf{M}_{(\mathbf{B})} = \mathbf{M} \otimes_{\mathbf{A}} \mathbf{B}\) 是秩为 \(n\) 的自由 B-模（利用定理 1 (d) 与命题 3）。证明其逆。

¶ 9. (a) 设 A 是一个环，\((f_i)_{i \in I}\) 是 A 的生成理想 A 的元素的有限族，M 是一个 A-模。设在每个 \(A_{f_i}\) -模 \(M_{f_i}\) 中给定了一个元素 \(z_i\) ，使得对 \(i \neq j\) ，\(z_i\) 与 \(z_j\) 在 \(M_{f_i f_j}\) 中的典范像相等。证明存在唯一的 \(z \in M\) ，使得对所有 \(i\) ，\(z\) 在 \(M_{f_i}\) 中的典范像等于 \(z_i\) 。（注意：对每个整数 \(k > 0\) ，存在 A 的元素的一个族 \((g_i)_{i \in I}\) ，使得 \(\sum_{i} g_i f_i^k = 1\) 。）

\begin{enumerate}
\def\labelenumi{(\alph{enumi})}
\setcounter{enumi}{1}
\item
  设 A 是一个环，P 是一个秩为 1 的投射 A-模。证明环 \(\operatorname{End}_{\mathbf{A}}(\mathbf{P})\) 典范地同构于 A（利用定理 3）。
\item
  设 A 是一个环，P 是一个秩为 n 的投射 A-模。对 P 的每个自同态 u，由 (b)，典范地对应于 \(\wedge^{n}u\) 的 A 的一个元素 \(\det(u)\) ，称为 u 的行列式，使得对 P 的两个自同态 u、v，有 \(\det(u \circ v) = (\det u)(\det v)\) 。元素 \(\chi_{u}(T) = \det(T.1 - u)\) （属于 A{[}T{]}）称为 u 的特征多项式（参看第 3 节命题 4）；证明 \(\chi_{u}\) 是次数为 n 的首一多项式，且 \(\chi_{u}(u) = 0\) （利用 (a) 及 §3 第 3 节定理 1）。由此推出：为使 u 是双射，必须且只需
\end{enumerate}

\(\det(u)\) 在 A 中可逆。若 A 是一个整环，\(\chi_u(T) = \prod_{i=1}^{n} (T - \alpha_i)\) 是 \(\chi_u\) 在 A 的分式域的一个代数闭扩张

中的线性因子分解，且 \(\chi_{q(u)}(T) = \prod_{i=1}^{l} (T - q(\alpha_i))\) 对每个多项式 \(q \in A[T]\) 都成立。类似地推广《代数》第 VII 章 § 5 第 6 节命题 14。

\begin{enumerate}
\def\labelenumi{\arabic{enumi}.}
\setcounter{enumi}{9}
\tightlist
\item
  设 A 是一个环，\(E(A)\) 是有限生成投射 A-模的同构类组成的集合；设 G 是 \(E(A)\) 的元素的形式线性组合组成的 Z-模，N 是 G 的由形如 \(\xi - \xi' - \xi''\) 的元素组成的子模，其中 \(\xi, \xi', \xi''\) 是三个投射 A-模 M、M′、M″ 的类，使得存在一个正合列
\end{enumerate}

\[
0 \rightarrow \mathrm{M} ^ {\prime} \rightarrow \mathrm{M} \rightarrow \mathrm{M} ^ {\prime \prime} \rightarrow 0
\]

（这等价于说，依《代数》第 II 章 §2 第 2 节，M 同构于直和 \(M' \oplus M''\) ）。设 \(K(A)\) 是商 Z-模 G/N，并且对每个有限生成投射 A-模，设 \(\kappa(M)\) （或 \(\kappa_{A}(M)\) ）是它在 \(K(A)\) 中的类。

\begin{enumerate}
\def\labelenumi{(\alph{enumi})}
\item
  证明在 \(K(\mathbf{A})\) 上存在唯一的交换环结构，其加法就是 Z-模 \(K(\mathbf{A})\) 的加法，其乘法满足：对两个有限生成投射 A-模 M、N，有 \(\kappa(\mathbf{M})\kappa(\mathbf{N}) = \kappa(\mathbf{M} \otimes_{\mathbf{A}} \mathbf{N})\) ；\(\kappa(\mathbf{A})\) 是 \(K(\mathbf{A})\) 的单位元。
\item
  对每个有限生成投射 A-模 M 及每个整数 \(p \geqslant 0\) ，设 \(\lambda^p\) 表示元素 \(\kappa(\bigwedge^p M)\) （在 \(K(A)\) 中）。证明
\end{enumerate}

\[
\lambda^ {p} (\mathbf {M} \oplus \mathbf {N}) = \sum_ {r + s = p} \lambda^ {r} (\mathbf {M}). \lambda^ {s} (\mathbf {N}).
\]

设 \(\lambda_{\mathrm{T}}(\mathbf{M})\) 表示元素 \(\sum_{p=0}^{\infty} \lambda^{p}(\mathbf{M}) \mathrm{T}^{p}\) ，它属于环 \(K(\mathrm{A})[[\mathrm{T}]]\) ，即环 \(K(\mathrm{A})\) 上一个未定元的形式幂级数环；证明

\[
\lambda_ {\mathbf {T}} (\mathbf {M} \oplus \mathbf {N}) = \lambda_ {\mathbf {T}} (\mathbf {M}) \lambda_ {\mathbf {T}} (\mathbf {N});
\]

由此导出一个同态，仍记作 \(x \mapsto \lambda_{\mathrm{T}}(x)\) ，它映 \(K(\mathrm{A})\) 到 \(K(\mathrm{A})[[\mathrm{T}]]\) 中常数项等于 1 的形式幂级数组成的乘法群。设 \(\lambda^{p}(x)\) 表示 \(T^{p}\) 在 \(\lambda_{\mathrm{T}}(x)\) 中的系数；那么

\[
\lambda^ {p} (x + y) = \sum_ {r + s = p} \lambda^ {r} (x) \lambda^ {s} (y)
\]

对所有 \(x, y\) （在 \(\mathbf{K}(\mathbf{A})\) 中）成立。

\begin{enumerate}
\def\labelenumi{(\alph{enumi})}
\setcounter{enumi}{2}
\item
  若 \(\mathbf{A}\) 是一个局部环，则环 \(K(\mathbf{A})\) 典范地等同于 \(\mathbf{Z}\) ，且 \(\lambda^p(n) = \binom{n}{p}\)。由此推出：若 \(\mathbf{A}\) 是 \(m\) 个局部环的直复合，则 \(K(\mathbf{A})\) 同构于 \(\mathbf{Z}^m\) 。
\item
  设 \(\phi: \mathbf{A} \to \mathbf{B}\) 是一个环同态。证明存在唯一的环同态 \(\phi': \mathbf{K}(\mathbf{A}) \to \mathbf{K}(\mathbf{B})\) ，使得 \(\phi'(\kappa_{\mathbf{A}}(\mathbf{M})) = \kappa_{\mathbf{B}}(\mathbf{M}_{(\mathbf{B})})\) 对每个有限生成投射 A-模 \(\mathbf{M}\) 成立；于是 \(\phi'(\lambda^p(x)) = \lambda^p(\phi'(x))\) 对所有 \(x \in \mathbf{K}(\mathbf{A})\) 成立；若 \(\psi: \mathbf{B} \to \mathbf{C}\) 是另一个环同态，则 \((\psi \circ \phi)' = \psi' \circ \phi'\) 。
\item
  设 \(\phi: A \to B\) 是一个环同态，\(B\) 关于它是一个有限生成投射 \(A\) -模；对每个有限生成投射 \(B\) -模 \(N\) ，\(\phi_{*}(N)\) 这时是一个有限生成投射 \(A\) -模；由此推出存在一个 \(Z\) -模同态 \(\phi_!: K(B) \to K(A)\) ，使得
\end{enumerate}

\[
\phi_! (\kappa_ {\mathrm{B}} (\mathrm{N})) = \kappa_ {\mathrm{A}} (\phi_ {*} (\mathrm{N})).
\]

证明：对所有 \(x \in K(A)\) 与 \(y \in K(B)\) ，有 \(\phi_!(y \cdot \phi'(x)) = \phi_!(y) \cdot x\) 。若 \(\psi: B \to C\) 是另一个使 \(C\) 为有限生成投射 \(B\) -模的环同态，则 \((\psi \circ \phi)_! = \phi_! \circ \psi_!\) 。

\begin{enumerate}
\def\labelenumi{\arabic{enumi}.}
\setcounter{enumi}{10}
\tightlist
\item
  设 K 是一个特征 \(\neq2\) 的域；考虑一个多项式 \(g(\mathbf{X})\in\mathbf{K}[\mathbf{X}]\) ，其次数为偶数且 \(\geqslant2\) ，它的所有根（在 K 的一个代数闭扩张中）都互不相同，并且满足 \(g(0)\neq0\) 。令 B = K{[}X{]}，
\end{enumerate}

\[
\mathbf {A} = \mathbf {B} [ \mathbf {Y} ] / (\mathbf {Y} ^ {2} - f (\mathbf {X}))
\]

其中 \(f(\mathbf{X}) = \mathbf{X}g(\mathbf{X})\) 。

\begin{enumerate}
\def\labelenumi{(\alph{enumi})}
\item
  设 \(\pmb{y}\) 是 \(\mathbf{Y}\) 在 \(\mathbf{A}\) 中的类。证明每个 \(a \in \mathbf{A}\) 都可以唯一地写成 \(a = \mathbf{P} + y\mathbf{Q}\) 的形式，其中 \(\mathbf{P}\) 与 \(\mathbf{Q}\) 属于 \(\mathbf{B}\) 。令 \(\bar{a} = \mathbf{P} - y\mathbf{Q}\) ，\(\mathrm{N}(a)=a\bar{a}=\mathrm{P}^{2}-f\mathrm{Q}^{2}\) 。证明关系 \(\mathrm{N}(a)=0\) 蕴含 a=0；由此推出 A 是一个整环。证明 A 的可逆元素就是 K 中 \(\neq0\) 的元素。
\item
  设 \(m\) 是理想 \(AX + Ay\) （在 \(A\) 中）；证明 \(m\) 是极大的，且理想 \(m_m\) （在 \(A_m\) 中）由 \(y\) 生成（注意 \(X = y^2 / g(X)\) ）。由此推出 \(m\) 是一个可逆 \(A\) -模（利用定理 1）。证明 \(m^2 = AX\) ；由此推出 \(m\) 不是 \(A\) 的主理想，并且 \(A\) 的分式域中可逆 \(A\) -模的类群不化为单位元。（若 \(m = At\) ，则 \(X = \lambda t^2\) ，其中 \(\lambda \in K\) ，从而 \(\lambda^{-1}X = P^2 + fQ^2\) ，其中 \(P \in B\) ，\(Q \in B\) ；证明这样的关系式不可能成立。）
\end{enumerate}

¶ 12. 设 A 是一个环，I 是一个 A-模，B 是以其底层加法群为 A ⊕ I 的唯一的环，其中 A 是 B 的子环，I 是平方为零的理想，且 I 经纯量限制得到的 A-模结构就是给定的结构（《代数》第 II 章 § 1 习题 7）。

\begin{enumerate}
\def\labelenumi{(\alph{enumi})}
\item
  假设 A 的每个不可逆元素都属于 I 的某个 \(\neq0\) 元素的零化子。证明 B 中每个不是零因子的元素都可逆，换言之，B 等于它的全分式环。
\item
  设 \((\mathfrak{m}_{\lambda})_{\lambda \in L}\) 是 A 的极大理想的一个族，使得 A 的每个不可逆理想都至少属于某个 \(\mathfrak{m}_{\lambda}\) 。证明：若取 I 为诸 A-模 \(\mathrm{A} / \mathfrak{m}_{\lambda}\) 的直和，则条件 (a) 得到满足。
\item
  由 (b) 推出：若存在一个非自由的秩 1 投射 A-模，则可以这样选取 I，使得 B 等于它的全分式环，但存在一个非自由的秩 1 投射 B-模（由于 B 是 B 的唯一的非退化子 B-模，这个模必然不是可逆的）。
\item
  假设 A 中存在一个极大理想 m，使得对所有 \(x \in m\) ，存在 A 的包含 x 的极大理想 \(m' \neq m\) 。若取 I 为诸 A/m′ 组成的直和 A-模，其中 \(m'\) 取遍 A 的异于 m 的极大理想组成的集合，则相应的环 \(B = A \oplus I\) 等于它的全分式环；若记 \(n = m \oplus I\) （它是 B 的一个极大理想），证明 \(B_n\) 同构于 \(A_m\) ；若 A 是一个整环（由假设必然不是域），则 \(B_n\) 是一个不同于其分式域的整环。若进一步 \(mA_m\) 是整环 \(A_m\) 的主理想，证明 B 的理想 \(mB = m \otimes_A B\) 是一个非自由的退化投射 B-模，秩为 1。（我们将在第 VII 章看到，存在这样的 Dedekind 整环 A：其中有极大理想 m，它的任何幂都不是主理想；对这样的环 A，上述所有假设都满足。）
\end{enumerate}

¶ 13. 设 A 是一个（交换）环，B 是一个 A-代数（不必交换）。假设 A-模 B 是忠实的，且由有限多个元素 \(b_i\) （ \(1 \leqslant i \leqslant n\) ）生成；回忆 \(B^0\) 表示 B 的反代数，并且 A-代数

B \(\otimes_{\mathbf{A}}\mathbf{B}^{0}\) 到 A-代数 \(\operatorname{End}_{\mathbf{A}}(\mathbf{B})\) 有一个典范同态，它把 \(b\otimes b'\) 映为自同态 \(z\mapsto bzb'\) 。

\begin{enumerate}
\def\labelenumi{(\alph{enumi})}
\tightlist
\item
  证明下列两个条件等价：
\end{enumerate}

\((\alpha)\) 诸 \(b_{i}\) 组成 A-模 B 的基，且典范同态 \(\mathbf{B} \otimes_{\mathbf{A}} \mathbf{B}^{0} \to \operatorname{End}_{\mathbf{A}}(\mathbf{B})\) 是双射。

(β) 矩阵 \(U = (u_{ij})\) （满足 \(u_{ij} = b_j b_i\) ）在 \(\mathbf{M}_n(\mathbf{B})\) 中可逆。

\begin{enumerate}
\def\labelenumi{(\alph{enumi})}
\setcounter{enumi}{1}
\tightlist
\item
  进一步假设 A 是以 m 为极大理想的局部环；设 \(\beta_{i}\) 表示 \(b_{i}\) 在 B/mB 中的典范像。证明：条件 ( \(\alpha\) ) 与 ( \(\beta\) ) 这时等价于下列每个条件：
\end{enumerate}

(γ) 诸 \(\beta_{i}\) 组成 A/m-模 B/mB 的基，且 B/mB 是一个单中心 A/m-代数（《代数》第 VIII 章 § 5 第 4 节）。

(δ) 矩阵 \((\beta_{j}\beta_{i})\) 在 \(\mathbf{M}_{n}(\mathrm{B}/\mathrm{mB})\) 中可逆。

（为证明 (δ) 蕴含 (β)，利用 mB 包含于 B 的 Jacobson 根这一事实，以及《代数》第 VIII 章 § 6 习题 5。）

\begin{enumerate}
\def\labelenumi{\arabic{enumi}.}
\setcounter{enumi}{13}
\tightlist
\item
  设 A 是一个（交换）环，B 是一个 A-代数，使得 A-模 B 是忠实平坦的且有限表示的。证明下列条件等价：
\end{enumerate}

(α) 对 A 的每个极大理想 m，A/m-代数 B/mB 都是单的且中心的。

(β) B 是一个投射 A-模，且若 \(B^{0}\) 表示 B 的反代数，则典范同态 \(B \otimes_{A} B^{0} \to \operatorname{End}_{A}(B)\) 是双射。

（归结到 A 是局部环的情形，并应用习题 13。）这时称 B 是 A 上的一个 Azumaya 代数。

¶ 15. 设 B 是一个 Azumaya A-代数（习题 14）。

\begin{enumerate}
\def\labelenumi{(\alph{enumi})}
\item
  证明 B 的中心与 B 的子环 A 相同，并且是 A-模 B 的一个直因子（参看第 I 章 § 2 习题 21 及第 II 章 § 5 习题 4）。
\item
  证明：对 B 的每个双边理想 \(\mathfrak{b}\) ，有 \(\mathfrak{b} = (\mathfrak{b} \cap A)B\) 。（首先注意 B 在 A-模 B 的每个自同态下稳定，并利用 B 到 A 的投射的存在性。）
\end{enumerate}

\begin{enumerate}
\def\labelenumi{\arabic{enumi}.}
\setcounter{enumi}{15}
\tightlist
\item
  \begin{enumerate}
  \def\labelenumii{(\alph{enumii})}
  \tightlist
  \item
    设 A 是一个（交换）环，A′ 是一个作为平坦 A-模的交换 A-代数，B 是一个 A-代数。证明：若 B 是一个 Azumaya 代数，则 \(B' = A' \otimes_{A} B\) 是一个 Azumaya A′-代数；若假设 \(A'\) 是一个忠实平坦 A-模，则其逆为真（参看《代数》第 II 章 §5 第 1 节命题 4 的推论，以及《交换代数》第 I 章 §2 第 9 节命题 10 与 §3 第 6 节命题 12）。
  \end{enumerate}
\end{enumerate}

\begin{enumerate}
\def\labelenumi{(\alph{enumi})}
\setcounter{enumi}{1}
\tightlist
\item
  证明 \(\mathbf{B} \otimes_{\mathbf{A}} \mathbf{C}\) （两个 Azumaya A-代数 \(\mathbf{B}, \mathbf{C}\) 的张量积）是一个 Azumaya A-代数（参看《代数》第 VIII 章 § 7 第 4 节定理 2 的推论 2）。
\end{enumerate}

¶ 17. 设 A 是一个环，P 是一个忠实的有限生成投射 A-模。

\begin{enumerate}
\def\labelenumi{(\alph{enumi})}
\item
  证明 A-代数 \(\mathbf{B} = \operatorname{End}_{\mathbf{A}}(\mathbf{P})\) 是一个 Azumaya 代数。（首先注意 A-模 B 同构于 \(\mathbf{P} \otimes_{\mathbf{A}} \mathbf{P}^*\) ，因而是有限表示的，然后利用 §2 第 7 节命题 19。）
\item
  证明 P 是一个左投射 B-模（利用定理 1 (d) 及 §3 习题 8 (c)）。
\item
  设 \(\mathbf{P}'\) 是另一个忠实的有限生成投射 A-模。证明 \(\mathbf{P} \otimes_{\mathbf{A}} \mathbf{P}'\) 是一个忠实的有限生成投射 A-模，且 \(\operatorname{End}_{\mathbf{A}}(\mathbf{P} \otimes_{\mathbf{A}} \mathbf{P}')\) 典范地等同于 \(\operatorname{End}_{\mathbf{A}}(\mathbf{P}) \otimes_{\mathbf{A}} \operatorname{End}_{\mathbf{A}}(\mathbf{P}')\) （利用定理 1 (e)、§ 2 第 7 节命题 18 与 19，以及 § 3 第 3 节定理 1）。考察 \(\mathbf{P}'\) 的秩为 1 的情形。
\end{enumerate}

¶ 18. 设 A 是一个环，P 是一个忠实的有限生成投射 A-模，B 是 A-代数 \(\operatorname{End}_{\mathbf{A}}(\mathbf{P})\) ，E 是一个左 B-模，F 是 A-模 \(\operatorname{Hom}_{\mathbf{B}}(\mathbf{P}, \mathbf{E})\) 。证明典范同态

\[
\beta \colon \mathrm{P} \otimes_ {\mathbf {A}} \mathrm{F} \rightarrow \mathrm{E}
\]

（它在《代数》第 VIII 章 § 1 第 4 节中定义）是双射。（按习题 17 (c) 的方式归结到 A 是局部环的情形；这时应用《代数》第 VIII 章 § 1 习题 9，并注意 \(\mathbf{F}^n\) 与 \(\operatorname{Hom}_{\mathbf{B}}(\mathbf{P}, \mathbf{E}^n)\) 是典范同构的 A-模。）

证明存在 F 的子 A-模组成的集合到 E 的子 B-模组成的集合上的一个严格递增双射（方法同上）。

¶ 19. 设 A 是一个环，B 与 C 是两个 A-代数，\(\phi: B \to C\) 是一个 A-代数同态。假设 B 是一个 Azumaya 代数（习题 14）。设 \(C'\) 是 C 的由与 \(\phi(B)\) 的元素交换的元素组成的子 A-代数。证明 C 同构于代数 \(B \otimes_A C'\) 。（把 C 视为 \(B \otimes_A B^0\) -模，并应用习题 18，同时注意 \(B \otimes_A B^0\) 同构于 \(\operatorname{End}_A(B)\) 。）

¶ 20. 设 A 是一个环，P、P′ 是两个忠实的有限生成投射 A-模。证明：若 θ: End\_A(P) → End\_A(P′) 是一个 A-代数同构，则存在一个秩为 1 的投射 A-模 F 和一个 A-同构 φ: P′ → P ⊗\_A F，使得 \(\theta(u) = \phi^{-1} \circ (u \otimes 1) \circ \phi\) 对所有 u ∈ End\_A(P) 成立。（对借助 θ 视为 (End\_A(P))-模的 E = P′ 应用习题 18；为证明 Hom\_B(P, P′) 是秩为 1 的投射 A-模，利用第 3 节定理 2 及《代数》第 VIII 章 § 5 习题 7 (c)。）

¶ 21. 设 A 是一个环，n 是一个 ≥0 的整数。

\begin{enumerate}
\def\labelenumi{(\alph{enumi})}
\tightlist
\item
  证明使得 \(F^{n}\) 同构于 \(A^{n}\) 的 A-模 F 都是秩为 1 的投射模，并且它们在 \(P(A)\) 中的类组成一个子群 \(P_{n}(A)\) ，其每个元素都被 n 零化。（为证明第一个断言，利用 §2 第 2 节的
\end{enumerate}

命题 5；为证明第二个断言，注意 \(\bigwedge^{n} \mathbf{F}^{n}\) 与 \(\otimes \mathbf{F}\) 是典范同构的。）

\begin{enumerate}
\def\labelenumi{(\alph{enumi})}
\setcounter{enumi}{1}
\tightlist
\item
  设 \(\operatorname{Aut}_n(A)\) 是矩阵 A-代数 \(\mathbf{M}_n(A)\) 的自同构组成的群，\(\operatorname{Int}_n(A)\) 是 \(\mathbf{M}_n(A)\) 的内自同构组成的子群。证明商群 \(\operatorname{Aut}_n(A)/\operatorname{Int}_n(A)\) 同构于 \(P_n(A)\) （取 \(P = P' = A^n\) 应用习题 20，或应用《代数》第 VIII 章 § 1 习题 9）。由此推出：若 \(A\) 是一个半局部环或主理想整环，则
\end{enumerate}

\[
\operatorname{Aut} _ {n} (\mathbf {A}) = \operatorname{Int} _ {n} (\mathbf {A}).
\]

* (c) 若 A 是一个 Dedekind 整环，证明 \(P_{n}(A)\) 与 \(P(A)\) 的由被 \(n_*\) 零化的元素组成的子群相同

¶ 22. 设 A 是一个环，M 是一个有限生成投射 A-模。为使 M 的一个有限表示子模 M′ 是 M 的直因子，必须且只需它是 M 的纯子模（利用定理 1 及 § 3 习题 15）。

\begin{enumerate}
\def\labelenumi{\arabic{enumi}.}
\setcounter{enumi}{22}
\tightlist
\item
  设 A 是一个环，\((f_i)_{i\in I}\) 是 A 的生成理想 A 的元素的一个有限族。
\end{enumerate}

\begin{enumerate}
\def\labelenumi{(\alph{enumi})}
\item
  设 \(\mathbf{M}\) 是一个 A-模。证明：若对所有 \(i\in \mathbf{I}\) ，\(\mathbf{M}_{f_i}\) 是一个有限生成 \(\mathbf{A}_{f_i}\) -模，则 \(\mathbf{M}\) 是一个有限生成 A-模。
\item
  设 \(\mathbf{B}\) 是一个（交换）A-代数。证明：若对所有 \(i\in \mathbf{I}\) ，\(\mathbf{B}_{f_i}\) 是一个有限生成 \(\mathbf{A}_{f_i}\) -代数，则 \(\mathbf{B}\) 是一个有限生成 A-代数。
\end{enumerate}

\begin{enumerate}
\def\labelenumi{\arabic{enumi}.}
\setcounter{enumi}{23}
\tightlist
\item
  设 \(\mathbf{A}\) 是一个使 \(\operatorname{Spec}(\mathbf{A})\) 连通的环。
\end{enumerate}

\begin{enumerate}
\def\labelenumi{(\alph{enumi})}
\item
  若 M 是一个秩为 n 的投射 A-模，则 M 的每个不化为 0 的直因子 N 都有一个 \(\leqslant n\) 的秩。
\item
  设 \(u\) 是有限秩投射 A-模 \(M\) 到有限秩投射 A-模 \(M'\) 的一个同态，使得 \(u(M)\) 是 \(M'\) 的一个直因子。证明 \(\operatorname{Ker}(u)\) 这时是 \(M\) 的直因子（参看习题 23），并且
\end{enumerate}

\[
\operatorname{rg} (\operatorname{Ker} (u)) + \operatorname{rg} (\operatorname{Im} (u)) = \operatorname{rg} (\mathbf {M}).
\]

此外，转置 \({}^t u\) 满足：\({}^t u(\mathbf{M}^*)\) 是 \(\mathbf{M}^*\) 的直因子，且 \(\operatorname{rg}({}^t u(\mathbf{M}^*)) = \operatorname{rg}(u(\mathbf{M}))\) （参看 § 3 习题 14 (a)）。

\begin{enumerate}
\def\labelenumi{\arabic{enumi}.}
\setcounter{enumi}{24}
\tightlist
\item
  设 A 是一个环，S 是 A 的一个乘性子集。A 的理想 a 称为 S-可逆的，如果存在 \(s \in S\) 及 A 的一个理想 b，使得 ab = As；A 的两个理想 a、b 称为 S-等价的，如果存在 S 中的 s、t，使得 sa = tb。证明 A 的 S-可逆理想的 S-等价类组成一个乘法群 \(C_{S}\) 。若 S 是饱和的（ \(\S 2\) 习题 1），且 S 中没有元素是 A 中的零因子，证明群 \(C_{S}\) 同构于 A 的全分式环 B（以及 \(S^{-1}A\) ）的可逆子 \(S^{-1}A\) -模的类组成的群。
\end{enumerate}

第 III 章（*）

